\chapter{Soundness, Completeness, Compactness, and Resolution}
\label{ch:prop-metatheory}

\section{Soundness}
\label{sec:prop-metatheory-01}

\begin{theorem}[Soundness of classical propositional natural deduction]
\label{thm:prop-nd-soundness}
For all \(\Gamma\subseteq\Fm_{\CPL}\) and
\(\varphi\in\Fm_{\CPL}\),
\[
 \Gamma\derives_{\ND_{\CPL}}\varphi
 \quad\Longrightarrow\quad
 \Gamma\entails\varphi.
\]
\end{theorem}
\begin{proof}
Induct on the construction of the derivation.  The induction statement is
stronger than the theorem: every valuation satisfying all open assumptions of
a subderivation satisfies its conclusion.

An assumption case is immediate.  The conjunction introduction and
elimination cases, disjunction introduction cases, implication elimination
case, negation elimination case, and bottom elimination case follow directly
from their recursive truth clauses.  For \(\lor E\), let \(v\) satisfy the
open assumptions of the conclusion.  It satisfies \(\varphi\lor\psi\), so it
satisfies \(\varphi\) or \(\psi\).  In the first case the induction hypothesis
for the left subsidiary derivation gives \(v\entails\chi\); in the second case
the right subsidiary derivation gives the same conclusion.

For \(\to I\), let \(v\) satisfy the open assumptions.  If
\(v\notentails\varphi\), then \(v\entails\varphi\to\psi\).  If
\(v\entails\varphi\), it satisfies every open assumption of the subsidiary
derivation, so the induction hypothesis gives \(v\entails\psi\).  Thus again
\(v\entails\varphi\to\psi\).  The \(\neg I\) case is analogous: if
\(v\entails\varphi\), the induction hypothesis would give
\(v\entails\bottom\), which is impossible; hence
\(v\entails\neg\varphi\).

For \(\mathsf{RAA}\), suppose \(v\) satisfies the open assumptions
and falsifies \(\varphi\).  Then it satisfies \(\neg\varphi\), so the
induction hypothesis for the subsidiary derivation gives
\(v\entails\bottom\), impossible.  Therefore \(v\entails\varphi\).  Every
rule preserves truth from open assumptions to conclusion, completing the
induction.
\end{proof}

\begin{theorem}[Soundness of the classical propositional sequent calculus]
\label{thm:prop-lk-soundness}
Every sequent derivable in \(\LK_{\CPL}\) is satisfied by every valuation.
\end{theorem}
\begin{proof}
Induct on the sequent derivation, showing that every derivable sequent is
satisfied by every valuation in the sense of the sequent-satisfaction
definition of \cref{sec:prop-deduction-02}.  Both forms of initial sequent
are satisfied by every valuation \(v\).  In an identity sequent
\(\Gamma,\varphi\Rightarrow\varphi,\Delta\), the formula \(\varphi\)
belongs to the support of the antecedent and to the support of the
succedent, so whichever truth value \(v\) gives \(\varphi\), either some
antecedent formula is false under \(v\) or some succedent formula is true
under \(v\).  In a \(\bottom\)-sequent
\(\Gamma,\bottom\Rightarrow\Delta\), the antecedent support contains
\(\bottom\), which is false under \(v\).

For the induction step, suppose toward a contradiction that some valuation
\(v\) does not satisfy the conclusion of a rule each of whose premises is
satisfied by every valuation.  By the sequent-satisfaction definition,
\(v\) makes every formula in the support of the conclusion's antecedent
true and every formula in the support of its succedent false.  For
weakening, the supports of the premise are included in those of the
conclusion, so \(v\) does not satisfy the premise.  For contraction,
premise and conclusion have the same supports, so exactly the same
valuations satisfy both.  For cut, \(v\) makes all formulas in
\(\Gamma,\Sigma\) true and all formulas in \(\Delta,\Pi\) false; because
\(v\) satisfies the left premise, \(\varphi\) is true under \(v\), and
because \(v\) satisfies the right premise, \(\varphi\) is false under
\(v\), which is impossible.

For each logical rule, use the truth clause for the principal connective to
produce a premise that \(v\) does not satisfy.  For example, if \(v\)
does not satisfy the conclusion of \(\land R\), then \(v\) makes every
formula in \(\Gamma\) true, every formula in \(\Delta\) false, and
\(\varphi\land\psi\) false; hence \(\varphi\) or \(\psi\) is false under
\(v\), so \(v\) does not satisfy the corresponding premise.  If \(v\)
does not satisfy the conclusion of \(\to L\), then \(v\) makes
\(\varphi\to\psi\) true.  If \(\varphi\) is false under \(v\), then
\(v\) does not satisfy the left premise; if \(\psi\) is true under
\(v\), then \(v\) does not satisfy the right premise.  The cases for
\(\neg L,\neg R,\land L,\lor L,\lor R,\to R\) follow by the same direct
application of their displayed truth clauses.  Thus every rule whose
premises are satisfied by every valuation has a conclusion satisfied by
every valuation, completing the induction.
\end{proof}

